\documentclass[10pt,a4paper]{amsart}
\usepackage[margin=19mm]{geometry}
\usepackage{amsmath,amssymb,mathtools}
\usepackage[scr=rsfs,cal=euler]{mathalfa}
\usepackage{xfrac}
\usepackage[unicode,hidelinks,bookmarks=false]{hyperref}
\newcommand{\ie}{i.e.\ }
\newcommand{\cf}{cf.\ }
\newcommand{\resp}{respectively\ }
\newcommand{\Subsection}{\subsection}
\providecommand{\nobreakdash}{\nobreak\hbox{-}}
\setlength{\emergencystretch}{2em}
\multlinegap0pt
\usepackage{xcolor}
\usepackage{amssymb}
\usepackage{algorithm}\usepackage{algpseudocode}
\usepackage[shortlabels]{enumitem}
\usepackage{tcolorbox}
\newtheorem{assumption}{Assumption}
\newtheorem{lemma}{Lemma}
\newtheorem{proposition}{Proposition}
\newcommand{\E}{\mathbb{E}}
\newcommand{\Farc}{\Phi}
\newcommand{\LS}{\operatorname{LS}}
\newcommand{\Oc}{\mathcal{O}}
\newcommand{\Proj}{\operatorname{Proj}}
\newcommand{\U}{\mathcal{U}}
\newcommand{\X}{\mathcal{X}}
\newcommand{\XPF}{{\mathcal{X}_{\rm PF}}}
\newcommand{\bx}{\mathbf{x}}
\newcommand{\hatFarc}{\hat{\Phi}}
\title{SURF: Steering the Scalarization Weight to Uniformly Traverse the Pareto Front}
\author{Liuyuan Jiang, Chentong Huang, Lisha Chen}
\date{Source: arXiv:2605.20619v1 (CC BY 4.0)}
\begin{document}
\maketitle

\noindent\emph{Source and licence.} SURF: Steering the Scalarization Weight to Uniformly Traverse the Pareto Front. Version arXiv:2605.20619v1 (CC BY 4.0), distributed by arXiv under the Creative Commons Attribution 4.0 International licence, http://creativecommons.org/licenses/by/4.0/, as recorded in the arXiv licence metadata for this exact version and shown on the abstract page https://arxiv.org/abs/2605.20619v1. Full licence notice: \texttt{licenses/arxiv-2605.20619-CC-BY-4.0.txt}.\par\smallskip

\noindent\textit{Editorial scope.} Counted unit: the article's proposition prop:bandit\_R\_and\_Farc\_err (source0.tex lines 1920--1932). The article's complete original proof of it is delivered with it (source0.tex lines 1933--2075, one \texttt{proof} environment of 143 lines). No mathematics is written, repaired or re-derived: the statement and the proof are the article's own lines, in the article's own notation, and no retained line is substituted or re-flowed.  Retained uncounted as the source-local context the proof needs: lines 342--345; lines 375--378; lines 390--393; lines 395--398.  Nothing was omitted from any retained span, so the package carries no omission record.  No retained line contains a citation, so the wrapper carries no bibliography and no bibliography entry.  No retained line contains a figure environment, an image-inclusion command or drawing code, so no image is delivered and the package declares no asset dependency.  Editorial formatting only, in addition: an AMS \texttt{amsart} preamble that restates the article's own declarations and macros used by the retained lines, namely the environments \texttt{assumption}, \texttt{lemma}, \texttt{proposition} on separate counters, together with the standard packages those lines need.  The retained environments print in this document as Assumption 1, Lemma 1, Proposition 1 in order of appearance, whereas the article prints the same items with the numbering of its own full document.  The complete native source of the article as distributed in the arXiv e-print for this exact version is bundled unchanged as \texttt{sources/201000/source0.tex}.
\begin{assumption}[Hidden-space regularity] 
\label{assump:hidden_regular} 
Assume there exists a bijection $g:\mathcal U\to\mathcal X$ such that $h_m(u)=f_m(g(u))$ for each $m\in[M]$. Each $f_m$ is twice continuously differentiable on an open set containing $\X$ and is $\mu$-strongly convex on $\mathcal X$. The hidden PS $\XPF =g(\U_{\rm PF})$ lies in a compact subset $\mathcal K\subset \mathcal X$ such that each $\nabla f_m$, is $l_{f,1}$-Lipschitz-continuous on $\mathcal K$. 
\end{assumption}

\begin{assumption}[Constraint Regularity]
\label{assump:LICQ}
Assume $\X=\{\bx:\mathbf c_1(\bx)\le 0,\mathbf c_2(\bx)= 0\}$, where $\mathbf c_1(\bx)$ is twice differentiable, Lipschitz-smooth and convex and $\mathbf c_2(\bx)$ is linear. For $\min_{\bx\in\X}\LS_f(\bx;w)$ with any $w\in[0,1]$, assume that the active constraints are fixed, satisfy the Linear Independence Constraint Qualification (LICQ), and the Lagrangian Hessian is Lipschitz-continuous along $\{\bx_w^*:w\in[0,1]\}$.
\end{assumption}

\begin{assumption}[Nonzero trade-off]
\label{assump:tangent_gap}
For each $w\in[0,1]$, let $J_w$ denote the Jacobian matrix of the active constraints of the scalarized problem $\min_{\bx\in\X}\LS_f(\bx;w)$ at $\bx_w^*$. Assume there exists some constant $c>0$ such that $\big\| \Proj_{\ker(J_w)}\big(\nabla f_1(\bx_w^*)-\nabla f_2(\bx_w^*)\big)\big\|\ge c$ for all $w\in[0,1]$.
\end{assumption}

\begin{lemma}[Bounded PF traversal speed]
\label{lem:v_bound_from_tangent_gap}
Suppose Assumptions~\ref{assump:hidden_regular},~\ref{assump:LICQ},~\ref{assump:tangent_gap} hold. Then, there exist constants $0<v_{\min}\le v_{\max}<\infty$ such that $v(w)\in [v_{\min},v_{\max}]$ for all $w\in[0,1]$. Consequently, $s(1)=\int_0^1 v(w) dw\in [v_{\min},v_{\max}]$, and $s(w)$ and $\Farc(w)$ are strictly increasing on $[0,1]$.
\end{lemma}

\begin{proposition}[Arc-length CDF estimation error in bandits]
\label{prop:bandit_R_and_Farc_err}
Consider MDP with singleton state space $|\mathcal{S}|=1$ and $|\mathcal{A}|=A\ge 2$ actions. Suppose the bounded speed condition in Lemma~\ref{lem:v_bound_from_tangent_gap} holds.
Given a budget of $T$ pulls, suppose that each pull of arm $a$ returns a bounded bi-objective reward vector
$[r_1(a),r_2(a)]$ with $\E[r_m(a)]=R_m(a)$ for $m\in\{1,2\}$.
We estimate $R_m(a)$ by the empirical mean $\hat R_m(a)$.
Then, with probability at least $1-\delta$, there is
\begin{align}
\|\hatFarc-\Farc\|_\infty
 = \Oc\left(s(1)^{-1}\|\hat R_m-R_m\|_\infty \right) = 
\Oc \left(s(1)^{-1}\sqrt{A\log(A/\delta)T^{-1}}\right).
\end{align}
\end{proposition}

\begin{proof}
We first control the reward estimation error.
Under a total budget of $T$ pulls, suppose each arm is sampled at least
$\lfloor T/A\rfloor$ times. Since each reward component lies in
$[-R_{\max},R_{\max}]$, Hoeffding's inequality yields that for any $\epsilon>0$ and
each $m\in\{1,2\}$,
$
\Pr \left(
\max_{a\in[A]} |\hat R_m(a)-R_m(a)| \ge \epsilon
\right)
\le
2A \exp \left(
-\frac{\lfloor T/A\rfloor \epsilon^2}{2R_{\max}^2}
\right).
$
Hence, by choosing $\epsilon_R
:=
2R_{\max}\sqrt{\frac{A\log(4A/\delta)}{2T}}$,
we obtain, with probability at least $1-\delta$,
\begin{align}
\|\hat R_m-R_m\|_\infty \le \epsilon_R,
\quad m\in\{1,2\}.
\label{eq:bandit_R_uniform_bound}
\end{align}

We now bound $\|\hatFarc-\Farc\|_\infty$. For any $w\in[0,1]$,
\begin{align*}
\big|\hatFarc(w)-\Farc(w)\big|
= & 
\left|\frac{\hat s(w)}{\hat s(1)}-\frac{s(w)}{s(1)}\right| \\
\leq &
\frac{|\hat s(w)-s(w)|}{s(1)}
+
\frac{s(w)}{s(1)}\frac{| \hat s(1)-s(1) |}{\hat s(1)}.
\end{align*}
Since $0\le s(w)\le s(1)$, we have
\begin{align*}
|\hat s(w)-s(w)|
\leq &
\int_0^w |\hat v_u(u)-v_u(u)| du
\le
\sup_{u\in[0,1]} |\hat v_u(u)-v_u(u)|,
\\
|\hat s(1)-s(1)|
\leq &
\int_0^1 |\hat v_u(u)-v_u(u)| du
\le
\sup_{u\in[0,1]} |\hat v_u(u)-v_u(u)|.
\end{align*}
Therefore,
\begin{align}
\big|\hatFarc(w)-\Farc(w)\big|
\le
\frac{1}{s(1)}\sup_{u\in[0,1]} |\hat v_u(u)-v_u(u)|
+
\frac{1}{\hat s(1)}\sup_{u\in[0,1]} |\hat v_u(u)-v_u(u)|.
\label{eq:Farc_diff_pre_event}
\end{align}

It remains to control the speed perturbation.
Since $|R_m(a)|\le R_{\max}$ for all $a\in[A]$ and $m\in\{1,2\}$,
the logits $\beta^{-1}(wR_1+(1-w)R_2)$ remain uniformly bounded in $\ell_\infty$
for all $w\in[0,1]$. Hence the softmax vector $\bx_w^*$ stays in a compact subset
of $\Delta_A^\circ$, uniformly over $w\in[0,1]$, and the map
$(R_1,R_2)\mapsto v(w)$ is uniformly Lipschitz-continuous over $w\in[0,1]$.
Therefore, there exists a constant $L_R<\infty$, depending only on
$(A,\beta,R_{\max})$, such that
\begin{align}
\sup_{w\in[0,1]} |\hat v(w)-v(w)|
\le
L_R\big(\|\hat R_1-R_1\|_\infty+\|\hat R_2-R_2\|_\infty\big).
\label{eq:bandit_speed_lipschitz}
\end{align}

On $\left\{ \sup_{u\in[0,1]} |\hat v_u(u)-v_u(u)| \le \frac{s(1)}{2} \right\}$, we have
$
\hat s(1)
\ge
s(1)-|\hat s(1)-s(1)|
\ge
s(1)-\sup_{u\in[0,1]} |\hat v_u(u)-v_u(u)|
\ge
\frac{s(1)}{2}.
$
Substituting this into \eqref{eq:Farc_diff_pre_event} gives, for all $w\in[0,1]$,
$
\big|\hatFarc(w)-\Farc(w)\big|
\le
\frac{1}{s(1)}\sup_{u} |\hat v_u(u)-v_u(u)|
+
\frac{2}{s(1)}\sup_{u} |\hat v_u(u)-v_u(u)|
=
\frac{3}{s(1)}\sup_{u\in[0,1]} |\hat v_u(u)-v_u(u)|.
$
Hence, on $\left\{ \sup_{u\in[0,1]} |\hat v_u(u)-v_u(u)| \le \frac{s(1)}{2} \right\}$,
\begin{align}
\|\hatFarc-\Farc\|_\infty
\le
\frac{3}{s(1)}\sup_{u\in[0,1]} |\hat v_u(u)-v_u(u)|.
\label{eq:Farc_bound_on_Ev}
\end{align}

Moreover, by \eqref{eq:bandit_speed_lipschitz}, the event $\left\{ \sup_{u\in[0,1]} |\hat v_u(u)-v_u(u)| \le \frac{s(1)}{2} \right\}$ is implied by
$
L_R\big(\|\hat R_1-R_1\|_\infty+\|\hat R_2-R_2\|_\infty\big)\le \frac{s(1)}{2}.
$
In particular, if
$
2L_R\epsilon_R\le \frac{s(1)}{2},
$
then \eqref{eq:bandit_R_uniform_bound} implies $\left\{ \sup_{u\in[0,1]} |\hat v_u(u)-v_u(u)| \le \frac{s(1)}{2} \right\}$.
Equivalently, it suffices that
$
\epsilon_R \le \frac{s(1)}{4L_R},
$
which holds whenever $T$ is large enough, namely
$
T
\gtrsim
\frac{A R_{\max}^2 L_R^2}{s(1)^2}\log \frac{A}{\delta}.
$

Finally, on the event \eqref{eq:bandit_R_uniform_bound}, combining
\eqref{eq:Farc_bound_on_Ev} and \eqref{eq:bandit_speed_lipschitz} yields
$
\|\hatFarc-\Farc\|_\infty
\le
\frac{3L_R}{s(1)}
\big(\|\hat R_1-R_1\|_\infty+\|\hat R_2-R_2\|_\infty\big)
\le
\frac{6L_R}{s(1)}\epsilon_R.
$
Therefore, with probability at least $1-\delta$,
$
\|\hatFarc-\Farc\|_\infty
=
O \left(
s(1)^{-1}\sqrt{\frac{A\log(A/\delta)}{T}}
\right),
$
where the hidden constant depends only on $(A,\beta,R_{\max})$ through $L_R$.
This completes the proof.
\end{proof}

\end{document}
